\documentclass[11pt]{article}
\usepackage[margin=0.8in, landscape]{geometry}
\usepackage{amsmath,amssymb,amsfonts}
\usepackage{tabularx}
\usepackage{booktabs}
\usepackage{multirow}
\usepackage{array}
\usepackage{float}
\usepackage{caption}
\usepackage[table]{xcolor}
\usepackage{colortbl}
\usepackage{hyperref}

\captionsetup{
  font=small,
  labelfont=bf,
  justification=centering,
  skip=6pt
}

\newcommand{\rev}[1]{#1}
\newcommand{\Cref}[1]{Table~#1}
\newcommand{\cref}[1]{Table~#1}

\pagestyle{empty}

\begin{document}

\vspace*{1cm}
\setcounter{table}{3}

\begin{table}[H]
\caption{Qualitative Comparison of Query Execution across Zero-Shot, Standard RAG, and Legalyze (Query ST-032).}
\label{tab:sample_comparison}
\centering
\small
\begin{tabularx}{\textwidth}{|l|>{\raggedright\arraybackslash\hsize=1.1\hsize}X|>{\raggedright\arraybackslash\hsize=1.1\hsize}X|>{\raggedright\arraybackslash\hsize=0.8\hsize}X|}
\hline
\textbf{Configuration} & \textbf{Retrieved Source Passage} & \textbf{Generated Response} & \textbf{Reviewer Evaluation} \\ \hline
\multicolumn{4}{|l|}{\textbf{Query (ST-032):} How does extortion under Section 383 PPC differ from theft under Section 378 PPC?} \\ \hline
\textbf{Zero-Shot} & 
\emph{No retrieval was used for this configuration.} & 
Extortion under Section 383 of the Pakistan Penal Code (PPC) differs from theft under Section 378 PPC. Extortion involves intentionally putting a person in fear of injury and thereby inducing them to deliver property. Theft, on the other hand, involves dishonestly taking movable property out of the possession of another person without their consent. & 
\textbf{N / Hallucinated} \newline (Unverified in context; no retrieved passage provided) \\ \hline
\textbf{Standard RAG} & 
\texttt{[Pakistan Penal Code, Section 379] Punishment for theft:} Whoever commits theft shall be punished with imprisonment of either description for a term which may extend to three years, or with fine, or with both. & 
Extortion under Section 383 PPC differs from theft under Section 378 PPC. According to Section 379, theft is punished with up to three years of imprisonment. Extortion under Section 383 PPC is when a person is put in fear of injury to deliver property, which is punished under Section 384 PPC with up to three years of imprisonment. & 
\textbf{N / Hallucinated} \newline (The response references Section 383 and 384, which are not present in retrieved Section 379) \\ \hline
\textbf{Legalyze} & 
\texttt{[Pakistan Penal Code, Section 378] Theft:} Whoever, intending to take dishonestly any moveable property out of the possession of any person without that person's consent, moves that property... \newline
\texttt{[Pakistan Penal Code, Section 383] Extortion:} Whoever intentionally puts any person in fear of any injury... and thereby dishonestly induces the person so put in fear to deliver... any property... & 
Under the Pakistan Penal Code (PPC), theft (Section 378) and extortion (Section 383) are distinguished by consent and coercion: 1) Theft involves taking moveable property out of someone's possession without their consent (Section 378). 2) Extortion involves inducing a person to deliver property by putting them in fear of injury (Section 383). & 
\textbf{Y / Traceable} \newline (Every claim is directly traceable to the retrieved source passages for Section 378 and 383) \\ \hline
\end{tabularx}
\end{table}

\end{document}
